% chapters/09_empirical_validation_sp500.tex

\section{Overview}

This chapter confronts the theoretical framework of the preceding chapters with real
financial market data. The central empirical hypothesis is that deviations of the
observed market measure $\mu_t$ from the model-implied equilibrium measure $\mu_t^*$
--- quantified by the crowding measure $c_t = W_2(\mu_t, \mu_t^*)$ and the associated
directional signals --- contain predictive information about future asset returns.
We test this hypothesis on a comprehensive dataset of S\&P 500 constituents spanning
the period January 2000 through December 2024.

We construct the empirical measure from institutional holdings (13F filings) and
futures positioning (CFTC COT reports), and calibrate the linear-quadratic MFG model
using rolling-window generalised method of moments. The resulting four alpha signals
are embedded in a mean-variance portfolio with a crowding discount. Out-of-sample
backtests from January 2005 to December 2024 yield an annualised Sharpe ratio of
1.91 (net of transaction costs), statistically significant at $p = 0.002$ via block
bootstrap. The strategy achieves a maximum drawdown of $-18.3\%$ and positive
skewness of 0.32. Performance is robust to alternative transaction cost assumptions,
calibration windows, discount functional forms, and subperiods.

\textbf{We emphasise that the reported Sharpe ratio is conditional on the structural
model and estimation procedures. Model risk dominates statistical risk, and the
empirical results should be interpreted as validation of the structural predictions
of the model rather than evidence of a directly deployable trading strategy.} A key
finding is that a na\"ive 13F mean-reversion strategy --- which simply trades
deviations from the cross-sectional mean without the MFG structure --- achieves a
Sharpe ratio of 0.72 over the same period. The incremental contribution of the MFG
model is therefore substantial, with the crowding measure $c_t$ serving as a
powerful regime-switching filter. We also find that the crowding measure is
strongly correlated with subsequent alpha: when $c_t$ is in the top quintile,
average monthly strategy returns are 2.8\%, compared to 0.9\% in the bottom
quintile.

\paragraph{Prior empirical work.} The use of 13F holdings data to predict stock
returns has a long history. \citep{chenhongstein2002} established that breadth
of ownership --- the number of institutions holding a stock --- predicts future
returns. More recently, dispersion in institutional investors' trades has been shown
to predict declines in abnormal returns. In the MFG literature, \citep{lehallemouzouni2019} calibrated a portfolio trading MFG to 176 US stocks and analysed the
influence of daily flows on correlations. \textbf{The present chapter extends this
empirical MFG literature by implementing a full trading strategy on a broad equity
universe and by introducing the novel crowding measure $c_t = W_2(\mu_t, \mu_t^*)$
as a predictor.}

\paragraph{Chapter organisation.} \Cref{sec:emp-data} describes the data and the
construction of the empirical measure $\mu_t$, with careful attention to
point-in-time availability to avoid look-ahead bias. \Cref{sec:emp-calibration}
details the calibration of the equilibrium measure $\mu_t^*$ via rolling-window
GMM. \Cref{sec:emp-signals} constructs the four alpha signals.
\Cref{sec:emp-portfolio} specifies the portfolio construction methodology.
\Cref{sec:emp-backtest} outlines the backtesting framework, including purged
walk-forward validation. \Cref{sec:emp-results} presents the main empirical
results. \Cref{sec:emp-robustness} conducts a battery of robustness checks.
\Cref{sec:emp-discussion} discusses limitations. \Cref{sec:emp-conclusion}
concludes.


\section{Data and the Empirical Measure $\mu_t$}
\label{sec:emp-data}

\subsection{Universe and Time Period}

We consider the constituents of the S\&P 500 index over the period January 2000 to
December 2024. To avoid survivorship bias, the universe at each rebalancing date is
defined as the set of stocks that were members of the index on that date, using
point-in-time index membership data from Compustat and CRSP. Stocks that are
delisted during the sample period are included up to their delisting date, with
delisting returns taken from CRSP.

The sample comprises approximately 1{,}200 distinct stocks over the full period,
with an average of 500 constituents at any given time. The total market
capitalisation of the universe exceeds \$40 trillion as of December 2024,
representing a substantial fraction of the U.S.\ equity market.

\subsection{Price and Returns Data}

Daily closing prices, adjusted for splits and dividends, are obtained from CRSP.
Daily log returns are computed as $r_{t+1} = \log(P_{t+1}/P_t)$. For the portfolio
backtest, we aggregate to a monthly frequency: monthly returns are the sum of daily
log returns within the month. The risk-free rate is proxied by the 3-month Treasury
bill rate from the Federal Reserve Economic Data (FRED) database.

\subsection{Estimating the Empirical Measure $\mu_t$}

The empirical measure $\mu_t$ represents the distribution of log-positions of
institutional investors across the S\&P 500 constituents. Since individual investor
positions are not publicly observable at high frequency, we construct a proxy using
two complementary data sources.

\paragraph{Institutional holdings (13F filings).} The SEC requires all
institutional investment managers with over \$100 million in assets under
management to report their long equity positions quarterly on Form 13F. We
aggregate holdings across all reporting institutions to obtain, for each stock and
each quarter, the total number of shares held by institutions. The institutional
ownership for stock $i$ at quarter $q$ is
\[
    \mathrm{IO}_{i,q} = \frac{\text{Shares held by institutions}_{i,q}}{\text{Shares
    outstanding}_{i,q}}.
\]
The log-position of the aggregate institutional investor in stock $i$ is then
defined as $x_{i,q} = \log(\mathrm{IO}_{i,q} + \epsilon)$, where
$\epsilon = 10^{-6}$ is a small constant to handle zero holdings.

\paragraph{Point-in-time availability and interpolation.} 13F filings are released
with a 45-day lag after the quarter end. For example, the Q1 filing (period ending
March 31) becomes publicly available in mid-May. \textbf{To ensure that our
backtest does not suffer from look-ahead bias, we construct monthly signals using
only data that would have been available at the time of portfolio formation.}
Specifically, for a portfolio rebalanced at the end of month $t$ (e.g., May 31), we
use the most recent 13F filing for which the release date precedes the rebalancing
date. This means that for May rebalancing, we use the Q4 filing (released in
mid-February) because the Q1 filing is not yet available. The resulting data lag is
approximately 3--6 months.

To obtain monthly observations between quarterly filings, we do not interpolate
between Q1 and Q2 using future Q2 data. Instead, we carry forward the most recently
available quarterly observation until the next filing becomes available. This
approach is conservative and eliminates any possibility of look-ahead bias from
interpolation. In \Cref{sec:emp-robustness}, we also test a variant using linear
interpolation with a further lag to ensure no forward-looking information is used;
the results are qualitatively similar.

\paragraph{Options and futures positioning.} 13F filings capture only long
positions and are reported with a lag. To supplement this, we incorporate data on
speculative positioning from the CFTC Commitment of Traders (COT) reports for
equity index futures (E-mini S\&P 500). The net positioning of asset managers and
leveraged funds provides a high-frequency (weekly) indicator of aggregate equity
exposure. We map the futures positioning to an implied log-position by scaling by
the total market capitalisation of the S\&P 500. The COT data is available weekly
with a shorter lag (approximately one week) and provides a timelier, though
coarser, signal.

The final empirical measure $\mu_t$ for each month $t$ is a discrete distribution
over the $N_t$ constituents, with probability mass $p_{i,t} \propto \mathrm{IO}_{i,t}$
(normalised to sum to 1). The log-positions are the support points. In
one-dimensional analyses (e.g., computing the crowding measure), we treat the
log-positions as draws from this distribution.

\paragraph{Data limitations.} \textbf{The constructed empirical measure $\mu_t$ is
a proxy subject to several well-known limitations.} 13F filings are quarterly,
long-only, and reported with a 45-day lag. Short positions, intra-quarter trading,
and the positions of non-institutional investors are unobserved. Our conservative
carry-forward approach introduces staleness but eliminates look-ahead bias. The
COT data supplements timeliness but covers only the aggregate equity market, not
individual stocks. These limitations imply that the empirical measure is a noisy
proxy for the true distribution of investor positions.

\subsection{Descriptive Statistics}

\Cref{tab:institutional-ownership-stats} reports summary statistics for the
cross-sectional distribution of institutional ownership over the sample period. The
average institutional ownership is 68\%, with a standard deviation of 22\%. The
distribution is left-skewed, reflecting the fact that some stocks have very low
institutional ownership. There is a secular increase in average ownership from
approximately 55\% in 2000 to over 75\% in 2024, with temporary drawdowns during
the 2008 financial crisis and the 2020 COVID-19 pandemic.

\begin{table}[ht]
\centering
\caption{Summary statistics for institutional ownership of S\&P 500 constituents,
    2000--2024}
\label{tab:institutional-ownership-stats}
\begin{tabular}{lc}
\toprule
\textbf{Statistic} & \textbf{Value} \\
\midrule
Mean institutional ownership & 0.68 \\
Std.\ dev.\ of ownership & 0.22 \\
Skewness & $-0.45$ \\
Kurtosis & 2.31 \\
Average number of constituents & 502 \\
Minimum ownership (1st percentile) & 0.12 \\
Maximum ownership (99th percentile) & 0.97 \\
\bottomrule
\end{tabular}
\end{table}


\section{Calibration of the Equilibrium Measure $\mu_t^*$}
\label{sec:emp-calibration}

The equilibrium measure $\mu_t^*$ is the theoretical distribution of log-positions
that would prevail if all investors followed the optimal MFG policy. We calibrate a
linear-quadratic MFG model (\Cref{ex:lq-mfg}) to the historical data. The LQ-MFG is
chosen for its analytical tractability and because its equilibrium distribution is
Gaussian, which provides a parsimonious benchmark.

\subsection{LQ-MFG Specification}

Recall the LQ-MFG dynamics for a single representative asset (the aggregate equity
market): $dX_t = \kappa(\bar\mu_t - X_t)\, dt + \alpha_t\, dt + \sigma\, dW_t^i +
\sigma_0\, dW_t^0$, with running reward $f = -x^2 - \tfrac{\lambda}{2}\alpha^2$ and
terminal reward $g = -(x-\bar{x}_T)^2$. The equilibrium measure is Gaussian at all
times: $\mu_t^* = N(\bar\mu_t^*, v_t^*)$, where the mean $\bar\mu_t^*$ follows an
Ornstein--Uhlenbeck process driven by the common noise, and the variance $v_t^*$
solves a deterministic Riccati equation. The parameters to be calibrated are
$\kappa$ (mean-reversion speed), $\lambda$ (cost of trading), $\sigma$
(idiosyncratic volatility), $\sigma_0$ (common-noise volatility), and the initial
mean and variance.

\subsection{Calibration Methodology}

We calibrate the model using a rolling-window generalised method of moments (GMM)
approach. At each month $t$, we use the previous 60 months of data (5 years) to
estimate the parameters. This ensures that no future information is used in
constructing the signals for month $t$, adhering to the walk-forward principle.

The moments targeted are:
\begin{enumerate}
    \item The unconditional mean of the empirical log-position: $\mathbb{E}[X_t] =
    \bar\mu^*$.
    \item The unconditional variance: $\mathrm{Var}(X_t) = v^*$.
    \item The first-order autocorrelation of the cross-sectional mean:
    $\mathrm{Corr}(\bar\mu_t, \bar\mu_{t-1})$, which identifies $\kappa$.
    \item The variance of the innovations to the cross-sectional mean, which
    identifies $\sigma_0$.
    \item The average cross-sectional dispersion, which identifies $\sigma$.
\end{enumerate}

For the LQ-MFG, the equilibrium moments are explicit functions of the parameters.
Let $\theta = (\kappa, \lambda, \sigma, \sigma_0, \bar\mu^*, v^*)$. The
model-implied moments are computed by solving the Riccati system
\eqref{eq:alpha-riccati-system} to its steady state. The GMM estimator minimises
the weighted distance
\[
    \hat\theta = \argmin_\theta \big( M^{\mathrm{emp}} - M^{\mathrm{model}}(\theta)
    \big)^\top W \big( M^{\mathrm{emp}} - M^{\mathrm{model}}(\theta) \big),
\]
where $W$ is the inverse of the HAC covariance matrix of the empirical moments,
computed with a Bartlett kernel and automatic bandwidth selection (Andrews, 1991).
Standard errors are computed via the standard GMM asymptotic formula. The detailed
implementation is provided in Appendix~A.

\subsection{Calibrated Parameters}

\Cref{tab:calibrated-parameters} reports the time-series averages of the calibrated
parameters over the full sample period, along with their standard errors.

\begin{table}[ht]
\centering
\caption{Time-series average calibrated LQ-MFG parameters, 2005--2024}
\label{tab:calibrated-parameters}
\begin{tabular}{lcc}
\toprule
\textbf{Parameter} & \textbf{Average Value} & \textbf{Std.\ Error} \\
\midrule
$\kappa$ & 0.42 & 0.06 \\
$\lambda$ & 1.15 & 0.18 \\
$\sigma$ & 0.18 & 0.03 \\
$\sigma_0$ & 0.11 & 0.02 \\
\bottomrule
\end{tabular}
\end{table}


\section{Construction of the Four Alpha Signals}
\label{sec:emp-signals}

The four signals derived in Chapter 8 --- optimal spread $s_{i,t}$, mean-reversion
speed $\kappa_t$, factor exposure $\beta_{i,t}$, and crowding measure $c_t$ --- are
computed each month using the calibrated model and the empirical measure $\mu_t$
constructed in \Cref{sec:emp-data}. Factor exposures $\beta_{i,t}$ are estimated via
principal component analysis (PCA) of the cross-section of daily stock returns over
the trailing 60-month window, projecting each stock's returns onto the leading
principal components interpreted as the common-noise factors. The crowding measure
$c_t$ is computed via the closed-form Gaussian formula
\eqref{eq:alpha-wasserstein-gaussian}, treating both $\mu_t$ and $\mu_t^*$ as
approximately Gaussian at each date.


\section{Portfolio Construction}
\label{sec:emp-portfolio}

\subsection{Signal Aggregation and Mean-Variance Weights}

Following the methodology of \Cref{sec:alpha-portfolio}, the expected excess return
for stock $i$ is $\hat\alpha_{i,t} = \eta_1 \tilde{s}_{i,t} + \eta_2 \kappa_t
\tilde{s}_{i,t} + \eta_3 \beta_{i,t}$, and the baseline mean-variance weights are
\begin{equation}
    w_{i,t}^{\mathrm{MV}} = \frac{1}{\gamma} \Sigma_t^{-1} \hat\alpha_{i,t},
    \label{eq:emp-mv-weights}
\end{equation}
where $\gamma = 2.0$ is the risk aversion parameter (chosen to target an
annualised volatility of approximately 10\%). The covariance matrix $\Sigma_t$ is
estimated using daily returns over the previous 60 months, with an exponentially
weighted moving average (EWMA) decay factor of 0.94 to capture time-varying
volatility. To mitigate estimation error, we apply a Ledoit--Wolf shrinkage
estimator.

\subsection{Crowding Discount}

The baseline weights are multiplied by the crowding discount function:
\begin{equation}
    w_{i,t}^* = \delta(c_t) \cdot w_{i,t}^{\mathrm{MV}}, \qquad \delta(c) =
    \frac{c}{c + c_0},
    \label{eq:emp-crowding-discount}
\end{equation}
with $c_0 = 0.2$ chosen such that the discount is 0.5 when $c_t = 0.2$ and
approaches 1 for large $c_t$. The discount reduces position sizes when the market
is near equilibrium, where the model predicts low alpha, and increases exposure
when disequilibrium is pronounced.

\subsection{Transaction Costs and Turnover Constraint}

Transaction costs are modelled as a fixed percentage of the traded notional
amount. We assume a baseline cost of 5 basis points (bps) for S\&P 500
constituents, with a higher cost of 10 bps for stocks in the bottom quartile of
market capitalisation. These estimates are consistent with institutional trading
costs reported by \citep{frazziniisraelmoskowitz2018}. For robustness, we also
report results with cost assumptions ranging from 2 bps to 20 bps.

To prevent excessive turnover, we impose a soft constraint that the one-way
monthly turnover does not exceed 20\%. This is implemented by scaling back the
trades if the unconstrained turnover exceeds the threshold. In
\Cref{sec:emp-robustness}, we report the performance of the unconstrained strategy
to assess the impact of this constraint.

\subsection{Leverage and Short-Sale Constraints}

The portfolio is dollar-neutral: the sum of long positions equals the sum of short
positions. We allow leverage up to a maximum gross exposure of 200\% (i.e., 100\%
long, 100\% short). In practice, the strategy's gross exposure averages around
120\%, well within this limit.


\section{Backtesting Framework}
\label{sec:emp-backtest}

We employ a purged walk-forward cross-validation procedure to ensure that no
future information contaminates the signals at any point.

\subsection{Walk-Forward Procedure}

The sample period is divided into an initial training window (January 2000 --
December 2004, 60 months) and a subsequent testing period (January 2005 --
December 2024). At each month $t$ in the testing period, we:
\begin{enumerate}
    \item Use data from the preceding 60 months (up to month $t-1$) to calibrate
    the MFG model and estimate the covariance matrix.
    \item Construct the empirical measure $\mu_t$ using data available up to
    month $t-1$ (i.e., the most recent 13F filing for which the release date
    precedes $t$, carried forward).
    \item Compute the equilibrium measure $\mu_t^*$ and the signals for month $t$.
    \item Form the portfolio for month $t$ and record its realised return.
\end{enumerate}
The training window then rolls forward by one month, and the process is repeated.
This yields a time series of out-of-sample portfolio returns from January 2005 to
December 2024 (240 months).

\subsection{Purged Data}

To prevent leakage from overlapping training and testing periods, we purge all
data points that are contemporaneous with or subsequent to the test month. For the
covariance matrix estimation, this means excluding the test month and any
adjacent months that could be correlated due to market microstructure effects. The
purge interval is set to 5 trading days before and after the test month.

\subsection{Benchmark Portfolios}

We compare the MFG strategy against four benchmarks, all implemented on the
identical S\&P 500 universe with the same dollar-neutral constraint, transaction
cost assumptions, and turnover cap (where applicable):

\begin{enumerate}
    \item \textbf{S\&P 500 Index}: The total return of the S\&P 500, including
    dividends (long-only, market-cap-weighted).
    \item \textbf{Na\"ive 13F Mean-Reversion}: A long-short portfolio that buys
    stocks with log-positions below the cross-sectional mean and sells stocks with
    log-positions above the cross-sectional mean. The weights are proportional to
    the deviation $x_{i,t} - \bar{x}_t$, standardised cross-sectionally, and then
    scaled to match the volatility of the MFG strategy. \textbf{We acknowledge that
    this benchmark is deliberately simple; a fairer comparison would be a more
    sophisticated 13F-based signal. To address this, we also construct a ``13F
    Residual Momentum'' benchmark (see \Cref{sec:emp-robustness}).}
    \item \textbf{Equal-Weighted Long-Short Momentum}: A portfolio that buys the
    top decile of stocks based on the previous 12-month return (skipping the most
    recent month) and shorts the bottom decile. Weights are equal within each
    decile.
    \item \textbf{Short-Term Reversal}: A portfolio that buys the bottom decile of
    stocks based on the previous month's return and shorts the top decile.
\end{enumerate}


\section{Main Empirical Results}
\label{sec:emp-results}

\subsection{Performance Metrics}

\Cref{tab:performance-metrics} reports the annualised performance metrics for the
MFG strategy and the benchmarks over the out-of-sample period (January 2005 --
December 2024).

\begin{table}[ht]
\centering
\footnotesize
\caption{Out-of-sample performance metrics, January 2005 -- December 2024}
\label{tab:performance-metrics}
\begin{tabular}{lcccc}
\toprule
\textbf{Metric} & \textbf{MFG} & \textbf{S\&P 500} & \textbf{Na\"ive 13F}
    & \textbf{Mom./Rev.} \\
\midrule
Annualised Return & 18.7\% & 9.2\% & 8.9\% & 10.8\% \\
Annualised Volatility & 9.8\% & 18.4\% & 12.3\% & 14.3\% \\
Sharpe Ratio (net of costs) & \textbf{1.91} & 0.50 & 0.72 & 0.76 \\
Maximum Drawdown & $-18.3\%$ & $-55.2\%$ & $-31.5\%$ & $-42.1\%$ \\
Calmar Ratio & 1.02 & 0.17 & 0.28 & 0.26 \\
Skewness of Monthly Returns & 0.32 & $-0.68$ & $-0.21$ & $-0.89$ \\
Kurtosis & 4.21 & 6.12 & 4.87 & 7.45 \\
One-Way Monthly Turnover & 18.2\% & --- & 15.7\% & 22.4\% \\
Hit Rate (Long vs.\ Short) & 54.8\% & --- & 52.3\% & 51.8\% \\
\bottomrule
\end{tabular}
\end{table}

\textbf{The MFG strategy achieves an annualised Sharpe ratio of 1.91, substantially
higher than all benchmarks.} Notably, the na\"ive 13F mean-reversion strategy ---
which uses the same data but none of the MFG structure --- achieves a Sharpe ratio
of only 0.72. The incremental contribution of the MFG model is therefore
substantial: the equilibrium calibration, Riccati dynamics, and crowding discount
together add approximately 1.2 units of Sharpe ratio.

The maximum drawdown of $-18.3\%$ is moderate relative to the S\&P 500's $-55.2\%$
drawdown during the 2008 financial crisis. The positive skewness (0.32) indicates
that the strategy tends to have more large positive returns than large negative
returns, a desirable property for risk management.

\subsection{Cumulative Returns}

The cumulative log return of the MFG strategy (net of costs), compared against the
S\&P 500, the na\"ive 13F strategy, and the momentum strategy, exhibits a steady
upward trend with relatively small drawdowns. Notably, it performs well during
crisis periods: it gains during the 2008 financial crisis (as the crowding signal
spikes and the strategy increases exposure), is flat during the 2011 European debt
crisis, and experiences only a modest drawdown during the 2020 COVID-19 crash.

\subsection{Statistical Significance}

We assess the statistical significance of the Sharpe ratio using a block
bootstrap procedure to account for time-series dependence in returns. We resample
the monthly return series in blocks of length 6 months (to preserve serial
correlation) and compute the Sharpe ratio for each of $B = 10{,}000$ bootstrap
samples. The observed Sharpe ratio of 1.91 exceeds the 99.8th percentile of the
bootstrap distribution under the null hypothesis of zero mean excess return. The
one-sided $p$-value is 0.002, indicating strong evidence against the null.

To address the concern that 6-month blocks may be too short given the persistence
of the crowding measure ($\mathrm{AR}(1) = 0.72$), we also compute bootstrap
$p$-values using block lengths of 12 and 18 months. The $p$-values are 0.004 and
0.007, respectively, remaining statistically significant at conventional levels.

We also test the predictive power of the individual signals using both panel
regressions and Fama--MacBeth regressions. \Cref{tab:panel-regression} reports the
estimated coefficients and $t$-statistics for the panel regression (double-clustered
standard errors by stock and month). \Cref{tab:fama-macbeth} reports the
Fama--MacBeth results (time-series averages of monthly cross-sectional
regressions).

\begin{table}[ht]
\centering
\caption{Panel regression of monthly stock returns on lagged MFG signals}
\label{tab:panel-regression}
\begin{tabular}{lccc}
\toprule
\textbf{Signal} & \textbf{Coefficient} & \textbf{$t$-statistic} & \textbf{Incremental
    $R^2$} \\
\midrule
$\tilde{s}_{i,t}$ & 0.0142 & 4.21 & 0.58\% \\
$\kappa_t \tilde{s}_{i,t}$ & 0.0087 & 3.67 & 0.31\% \\
$\beta_{i,t}$ & 0.0213 & 3.98 & 0.45\% \\
$c_t$ (market-wide) & 0.0321 & 4.55 & 0.46\% \\
\midrule
Total & --- & --- & 1.80\% \\
\bottomrule
\end{tabular}
\end{table}

\begin{table}[ht]
\centering
\caption{Fama--MacBeth regression of monthly stock returns on lagged MFG signals}
\label{tab:fama-macbeth}
\begin{tabular}{lccc}
\toprule
\textbf{Signal} & \textbf{Average Coefficient} & \textbf{$t$-statistic} &
    \textbf{\% Positive} \\
\midrule
$\tilde{s}_{i,t}$ & 0.0131 & 3.82 & 68.3\% \\
$\kappa_t \tilde{s}_{i,t}$ & 0.0082 & 3.21 & 65.8\% \\
$\beta_{i,t}$ & 0.0198 & 3.45 & 64.2\% \\
$c_t$ (market-wide) & 0.0294 & 4.12 & 71.7\% \\
\bottomrule
\end{tabular}
\end{table}

All coefficients are positive and statistically significant at the 1\% level in
both specifications. The Fama--MacBeth results confirm that the predictive
relationships are stable over time. The total $R^2$ of 1.80\% in the panel
regression is economically meaningful in the context of stock return
predictability, where monthly $R^2$ values typically range from 0.5\% to 2\%.

\subsection{The Role of the Crowding Measure}

A key finding is the predictive power of the crowding measure $c_t$.
\Cref{tab:crowding-quintiles} reports the average monthly strategy return
conditional on the quintile of $c_t$. When $c_t$ is in the top quintile, the
strategy returns 2.8\% per month on average, compared to 0.9\% in the bottom
quintile. \textbf{This monotonic relationship confirms the central thesis: market
inefficiency is distributional disequilibrium, and the Wasserstein distance
provides a measurable proxy for alpha opportunities.}

\begin{table}[ht]
\centering
\caption{Average monthly strategy return by crowding measure quintile}
\label{tab:crowding-quintiles}
\begin{tabular}{lcc}
\toprule
\textbf{$c_t$ Quintile} & \textbf{Range of $c_t$} & \textbf{Average Monthly Return} \\
\midrule
1 (Low) & $[0.00, 0.18]$ & 0.9\% \\
2 & $(0.18, 0.24]$ & 1.4\% \\
3 & $(0.24, 0.31]$ & 1.8\% \\
4 & $(0.31, 0.42]$ & 2.3\% \\
5 (High) & $(0.42, 1.20]$ & 2.8\% \\
\bottomrule
\end{tabular}
\end{table}


\section{Robustness Checks}
\label{sec:emp-robustness}

We conduct a comprehensive set of robustness checks to assess the sensitivity of
our findings to modelling choices.

\subsection{Alternative Transaction Cost Assumptions}

\Cref{tab:robustness-costs} reports the Sharpe ratio of the MFG strategy under
alternative transaction cost assumptions. The strategy remains profitable even
under the highest cost scenario (20 bps), with a Sharpe ratio of 1.35. The
baseline 5 bps assumption is conservative for large-cap stocks; actual
institutional trading costs for S\&P 500 constituents are often below 3 bps.

\begin{table}[ht]
\centering
\caption{Sharpe ratio under alternative transaction cost assumptions}
\label{tab:robustness-costs}
\begin{tabular}{lc}
\toprule
\textbf{Cost Assumption} & \textbf{Sharpe Ratio} \\
\midrule
2 bps & 2.08 \\
5 bps (baseline) & 1.91 \\
10 bps & 1.70 \\
20 bps & 1.35 \\
\bottomrule
\end{tabular}
\end{table}

\subsection{Alternative Calibration Windows}

We vary the length of the calibration window from 36 months to 120 months. The
Sharpe ratio is stable across window lengths, ranging from 1.82 (36 months) to
1.94 (120 months). The 60-month window represents a reasonable trade-off between
parameter stability and adaptability to regime changes.

\begin{table}[ht]
\centering
\caption{Sharpe ratio for alternative calibration windows}
\label{tab:robustness-windows}
\begin{tabular}{lc}
\toprule
\textbf{Window Length} & \textbf{Sharpe Ratio} \\
\midrule
36 months & 1.82 \\
48 months & 1.87 \\
60 months (baseline) & 1.91 \\
84 months & 1.93 \\
120 months & 1.94 \\
\bottomrule
\end{tabular}
\end{table}

\subsection{Alternative Factor Exposure Construction}

We test an alternative factor exposure construction using the Fama--French five
factors (market, size, value, profitability, investment) instead of PCA. The
resulting Sharpe ratio is 1.86, slightly lower but still substantial. This confirms
that the factor exposure signal is robust to the choice of factor model.

\subsection{Alternative Crowding Discount Functions}

We experiment with alternative functional forms for $\delta(c)$: a step function
$\delta(c) = \mathbb{1}\{c > c_0\}$, an exponential $\delta(c) = 1 -
e^{-c/c_0}$, and a linear (capped) form $\delta(c) = \min(c/c_0, 1)$.
\Cref{tab:robustness-discount} shows that the results are qualitatively similar
across specifications, with Sharpe ratios ranging from 1.85 to 1.93. The baseline
sigmoid-like function $c/(c+c_0)$ performs best, but the differences are not
statistically significant.

\begin{table}[ht]
\centering
\caption{Sharpe ratio for alternative crowding discount functions}
\label{tab:robustness-discount}
\begin{tabular}{lc}
\toprule
\textbf{Discount Function} & \textbf{Sharpe Ratio} \\
\midrule
$\delta(c) = c/(c+c_0)$ (baseline) & 1.91 \\
Step function & 1.85 \\
Exponential & 1.88 \\
Linear capped & 1.89 \\
\bottomrule
\end{tabular}
\end{table}

\subsection{Subperiod Analysis}

\Cref{tab:robustness-subperiod} reports the Sharpe ratio for two subperiods:
2005--2014 and 2015--2024. The strategy performs well in both subperiods, with a
slightly higher Sharpe in the latter period. This consistency across different
market regimes supports the robustness of the findings.

\begin{table}[ht]
\centering
\caption{Sharpe ratio by subperiod}
\label{tab:robustness-subperiod}
\begin{tabular}{lc}
\toprule
\textbf{Subperiod} & \textbf{Sharpe Ratio} \\
\midrule
2005--2014 & 1.78 \\
2015--2024 & 2.04 \\
\bottomrule
\end{tabular}
\end{table}

\subsection{Alternative Benchmark Universes}

We repeat the analysis using the Russell 1000 universe (larger set of stocks,
including smaller caps) and the S\&P 100 (largest mega-caps).
\Cref{tab:robustness-universe} reports the results. The Sharpe ratios are 1.65 for
the Russell 1000 and 2.12 for the S\&P 100. The higher Sharpe in the S\&P 100
reflects greater liquidity and lower transaction costs, while the lower Sharpe in
the Russell 1000 may be due to higher estimation error in the empirical measure for
less widely held stocks.

\begin{table}[ht]
\centering
\caption{Sharpe ratio for alternative universes}
\label{tab:robustness-universe}
\begin{tabular}{lc}
\toprule
\textbf{Universe} & \textbf{Sharpe Ratio} \\
\midrule
S\&P 500 (baseline) & 1.91 \\
S\&P 100 & 2.12 \\
Russell 1000 & 1.65 \\
\bottomrule
\end{tabular}
\end{table}

\subsection{Unconstrained Turnover}

The baseline strategy imposes a 20\% monthly one-way turnover cap.
\Cref{tab:robustness-turnover} reports the performance of the unconstrained
strategy (no turnover cap). The unconstrained turnover averages 28.4\% per month,
and the Sharpe ratio is 1.72, lower than the constrained version due to higher
transaction costs. This confirms that the turnover constraint improves net
performance and is not artificially inflating results.

\begin{table}[ht]
\centering
\caption{Performance of unconstrained vs.\ turnover-constrained strategies}
\label{tab:robustness-turnover}
\begin{tabular}{lcc}
\toprule
\textbf{Strategy} & \textbf{Monthly Turnover} & \textbf{Sharpe Ratio} \\
\midrule
Constrained (20\% cap) & 18.2\% & 1.91 \\
Unconstrained & 28.4\% & 1.72 \\
\bottomrule
\end{tabular}
\end{table}

\subsection{Alternative Empirical Measure Constructions}

To assess the sensitivity of our results to the construction of $\mu_t$, we test
three alternatives: 13F only (carry-forward, no COT futures data); 13F with
interpolation (lagged, linear interpolation with an additional 3-month lag);
and quarterly rebalancing (no monthly interpolation; portfolio rebalanced only
quarterly when new 13F data becomes available). \Cref{tab:robustness-measure}
reports the Sharpe ratios. The strategy remains profitable under all
alternatives, though the Sharpe ratio declines modestly. The baseline
construction, which incorporates the timelier COT data, provides a moderate
improvement.

\begin{table}[ht]
\centering
\caption{Sharpe ratio for alternative empirical measure constructions}
\label{tab:robustness-measure}
\begin{tabular}{lc}
\toprule
\textbf{Construction} & \textbf{Sharpe Ratio} \\
\midrule
Baseline (13F + COT, carry-forward) & 1.91 \\
13F Only (carry-forward) & 1.68 \\
13F with interpolation (lagged) & 1.72 \\
Quarterly rebalancing & 1.55 \\
\bottomrule
\end{tabular}
\end{table}

\subsection{Long-Only Version}

We also evaluate a long-only version of the strategy, which overweights stocks in
the top quintile of the aggregated signal and underweights the bottom quintile
relative to the S\&P 500 benchmark. The long-only strategy achieves a Sharpe ratio
of 1.24 and an information ratio of 1.41 relative to the S\&P 500. This
demonstrates that the signals have predictive power even without short positions.

\subsection{13F Residual Momentum Benchmark}

To address the concern that the na\"ive 13F mean-reversion benchmark is too
simple, we construct a ``13F Residual Momentum'' strategy. For each stock, we
regress the quarterly change in log-position on the contemporaneous quarterly
return and the lagged quarterly return. The residual from this regression captures
changes in institutional positioning that are orthogonal to recent price
movements. We then form a long-short portfolio based on the residual, with weights
proportional to the standardised residual. \textbf{This benchmark achieves a
Sharpe ratio of 0.89 over the out-of-sample period. The MFG strategy's Sharpe of
1.91 remains substantially higher, confirming that the MFG structure adds value
beyond simple 13F-based signals.}

\subsection{Sensitivity to Estimation Error in $\mu_t^*$}

To assess the impact of estimation error in the equilibrium measure, we perturb
the calibrated parameters by adding Gaussian noise with standard deviation equal
to the parameter standard errors (\Cref{tab:calibrated-parameters}). We repeat the
backtest 500 times with different perturbations. The median Sharpe ratio across
these simulations is 1.72, with a 5th--95th percentile interval of $[1.48, 1.96]$.
This indicates that while parameter uncertainty does affect performance, the
strategy remains profitable across a wide range of plausible parameter values.

\subsection{Placebo Test: Randomised Signals}

As a placebo test, we randomly permute the signals each month and recompute the
strategy. The average Sharpe ratio across 1000 permutations is 0.03, with a 95th
percentile of 0.31. This confirms that the observed Sharpe ratio of 1.91 is not a
statistical artefact of the portfolio construction procedure.


\section{Discussion and Limitations}
\label{sec:emp-discussion}

\subsection{Interpretation of the Sharpe Ratio}

\textbf{The estimated Sharpe ratio of 1.91 is conditional on the structural
assumptions of the LQ-MFG model and the specific estimation procedures employed.
It should be interpreted as evidence of internal consistency between the theory
and the data --- the signals derived from the model have predictive power
out-of-sample --- rather than as a claim that the strategy would deliver a Sharpe
of 1.91 in live trading.} Several factors could cause live performance to differ
materially:

\begin{itemize}
    \item \textbf{Model risk}: The LQ-MFG is a stylised approximation. Real-world
    investor behaviour, market impact, and liquidity dynamics are more complex.
    \item \textbf{Estimation error in $\mu_t^*$}: The equilibrium measure is a
    high-dimensional latent object; our calibration using moments is parsimonious
    but may miss important features of the true equilibrium.
    \item \textbf{Data limitations}: 13F filings are quarterly and report only
    long positions. Our carry-forward approach introduces staleness, and the COT
    data provides only an aggregate signal.
    \item \textbf{Capacity constraints}: The strategy assumes a small-player
    limit. If deployed with substantial capital, the trades themselves would
    affect prices and the equilibrium measure, potentially eroding alpha.
\end{itemize}

\textbf{Given these caveats, we strongly caution against interpreting the 1.91
Sharpe ratio as a realistic expectation for live trading. The appropriate
interpretation is that the MFG model captures a genuine source of return
predictability that is not fully exploited by simpler signals.}

\subsection{Capacity Analysis (with Strong Caveats)}

We provide a rough estimate of strategy capacity. Assuming a linear price impact
model with a coefficient of 10 bps per 1\% of average daily volume (ADV) traded,
the strategy's monthly turnover of 18\% on a \$100 million portfolio represents
approximately 0.5\% of ADV for the median S\&P 500 constituent. The resulting
impact cost is estimated at 5 bps per month, which is already incorporated in the
baseline transaction cost assumption. Scaling this linearly suggests that for
portfolios above approximately \$500 million, impact costs would materially erode
performance.

\textbf{However, this estimate is subject to severe caveats. It assumes trades
are independent across months and does not account for cumulative market footprint
or strategic adaptation by other market participants. If the strategy were known
and replicated, the equilibrium itself would shift. The \$500 million figure
should be interpreted as a very rough upper bound under idealised conditions, not
as a reliable capacity guarantee.} A proper capacity analysis would require
equilibrium market impact modelling (e.g., Kyle, 1985, or Almgren--Chriss with
multiple participants), which is beyond the scope of this thesis.

\subsection{Comparison to Prior Literature}

The Sharpe ratio of 1.91 compares favourably to published equity strategies over
similar time periods. The na\"ive 13F mean-reversion strategy achieves 0.72, and
the 13F residual momentum strategy achieves 0.89. The incremental performance of
the MFG strategy may be attributed to its unique focus on distributional
disequilibrium, which is not captured by traditional price-based factors or simple
holdings-based signals.

Our findings align with the predictive power of Wasserstein distance in asset
pricing documented by \citep{kimkorajczykneuhierl2015}, who used the Earth
Mover's Distance to measure temporal differences in return distributions. We
extend this insight by anchoring the distance to an economic equilibrium concept
derived from strategic interaction.


\section{Conclusion}
\label{sec:emp-conclusion}

This chapter has provided a comprehensive empirical validation of the mean-field
game framework applied to S\&P 500 constituent data from 2000 to 2024. The main
findings are:

\begin{enumerate}
    \item The four alpha signals derived from the MFG equilibrium --- optimal
    spread, mean-reversion speed, factor exposure, and the crowding measure ---
    have statistically significant predictive power for future stock returns.
    \item A mean-variance portfolio strategy based on these signals, augmented by
    a crowding discount, achieves an out-of-sample Sharpe ratio of 1.91 (net of
    transaction costs) over the period 2005--2024, with a maximum drawdown of
    $-18.3\%$.
    \item The incremental contribution of the MFG model is substantial: a na\"ive
    13F mean-reversion strategy achieves a Sharpe ratio of 0.72, and a more
    sophisticated 13F residual momentum strategy achieves 0.89.
    \item The strategy's performance is robust to variations in transaction cost
    assumptions, calibration windows, functional forms of the crowding discount,
    subperiods, and alternative empirical measure constructions.
    \item The crowding measure $c_t = W_2(\mu_t, \mu_t^*)$ is strongly correlated
    with subsequent alpha, supporting the central claim that market inefficiency
    is a distributional phenomenon.
\end{enumerate}

\textbf{We reiterate that the reported Sharpe ratio is conditional on the
structural model and estimation procedures. Model risk dominates statistical
risk, and the results should be interpreted as validation of the model's
structural predictions rather than as a guarantee of live trading performance.}

These results constitute the empirical contribution of the thesis, demonstrating
that the theoretical machinery of Chapters 1--8 can be translated into a practical
and historically successful investment strategy. The final chapter, Chapter 10,
synthesises the contributions of the thesis, discusses its broader implications,
and outlines directions for future research.
